% SPDX-License-Identifier: Apache-2.0
\section{A Cast}
\label{sec:x-cast}

\code{x as T} lowered as \code{x}: the cast was dropped. The macro was
made to log every cast it dropped, and the tree held 132, every one
widening a value of at most 32 bits or casting a constant, so no
netlist was wrong. It would have been the first time one narrowed. A
cast to an unsigned integer, \code{u8} to \code{u128} or
\code{usize}, now keeps the low bits of the value, as Rust's does:
the netlist has the value itself where it is no wider, which leaves
all 132 as they were, and its low bits where it is wider. A cast to
anything else is refused with a message naming \code{resize} and
\code{sext} (issue 496).

The unit is a ring of 256 bytes whose twelve-bit pointers count on
past its end, each addressing a word as
\code{p.raw() as u8 as usize}. The run asserts that 300 words come
back in the order they went in, and the netlist is checked against
the run under both simulators. With the cast dropped as it used to
be, the netlist addressed word 256 where the run addressed word 0, and
nvc stopped on an index out of range; Verilator did not, which is how
the old behaviour could have gone unseen.

\begin{figure}[htbp]
\centering
\input{cast_timing.tex}
\caption{The ring as the pointers pass the end: \code{wp} and
\code{rp} go on counting, and the words read are still the ones
written eight earlier.}
\label{fig:x-cast}
\end{figure}

\lstinputlisting[caption={\code{ex\_cast.rs}. Run with \code{bazel run //lib/examples:ex\_cast}.},label={lst:x-cast}]{lib/examples/ex_cast.rs}

\lstinputlisting[style=ascii,caption={What \code{ex\_cast} printed when the build ran it, then the lowering.},label={lst:x-cast-out}]{out_cast.txt}
